\uuid{Avcs}
\titre{ Facteur $\frac{1}{2m}$ dans la MSE : impact sur la convergence }
\niveau{L3}
\module{Analyse de données}
\chapitre{Réseaux de neurones}
\sousChapitre{Fonctions de coût, descente de gradient}
\theme{MSE, facteur de normalisation, minimum, convergence, learning rate}
\auteur{Maxime Nguyen}
\datecreate{2026-03-19}
\organisation{AMSCC}
\difficulte{3}

\contenu{
\texte{
  On considère deux variantes du coût MSE :
  \[
    J_1 = \frac{1}{m} \sum_{i=1}^m (\hat{y}^{(i)} - y^{(i)})^2, \qquad
    J_2 = \frac{1}{2m} \sum_{i=1}^m (\hat{y}^{(i)} - y^{(i)})^2.
  \]
  Un code utilise $J_1$ au lieu de $J_2$.
}
\begin{enumerate}

\item
  \question{La descente de gradient avec $J_1$ converge-t-elle vers le même minimum qu'avec $J_2$ ?}
  \reponse{
    Oui. $J_1 = 2 J_2$, donc les deux fonctions ont exactement les mêmes minimiseurs :
    annuler le gradient de l'une annule celui de l'autre. Le minimum est identique.
  }

\item
  \question{Les poids obtenus seront-ils identiques ?}
  \reponse{
    Oui, à convergence. Le minimum global est le même pour $J_1$ et $J_2$,
    donc les poids optimaux $w^*$ sont identiques.
  }

\item
  \question{Le nombre d'itérations nécessaires changera-t-il ?}
  \reponse{
    Oui, en pratique. Le gradient de $J_1$ est exactement le double de celui de $J_2$ :
    \[
      \nabla J_1 = 2\,\nabla J_2.
    \]
    Avec le même taux d'apprentissage $\alpha$, les mises à jour sont deux fois plus grandes avec $J_1$.
    Cela peut accélérer la convergence si $\alpha$ est bien choisi, ou provoquer des oscillations
    si $\alpha$ est trop grand. Le facteur $\frac{1}{2}$ dans $J_2$ simplifie simplement
    le calcul analytique du gradient (le $2$ du carré se simplifie) sans changer le problème d'optimisation.
  }

\end{enumerate}
}
